% Chapter 54: How Atoms Form Molecules and Solids
\chapter{How Atoms Form Molecules and Solids}

Chapter~53 traced a hundred elemental personalities to outer electron counts. But isolated atoms are only the world's parts bin. Almost everything around you, water, salt, plastic, steel, glass, muscle, chips, consists of molecular or solid assemblies. Different arrangements of identical atoms make radically different things. What happens as atoms approach that sometimes binds them tightly, sometimes leaves them indifferent, and sometimes builds conducting, brittle, or ductile solids? The answer leads to your phone chip's heart.

\step{A Counterintuitive Opening}

Pencil lead and diamond rings are the same thing.

Literally false, since pencil cores contain clay and diamonds impurities, but both pure substances are carbon. Diamond and graphite have identical carbon atoms and electron counts. Yet diamond is nature's hardest material, transparent and insulating; graphite sheds writing flakes, is black, and conducts well enough for electrodes. Why does one cut glass while the other lubricates? Why does one block current and the other admit it?

Another pair rules out coincidence. Salt and iron are crystalline solids, yet a salt grain shatters against a table while iron wire bends, flattens, and draws thin without breaking. One orderly atomic assembly is biscuit-brittle, another clay-tough. Stranger, solid salt insulates but saltwater conducts, while copper conducts both molten and solid.

Finally your phone. A fingernail-sized silicon chip contains tens of billions of switches, each neighboring regions of slightly different composition in one crystal. Without moving parts or vacuum tubes, how does inert-looking gray silicon decide whether current passes, or amplify weak signals into headphone sound?

\step{Make a Guess First}

Write three predictions first.

First: hydrogen atoms bind into molecules, helium atoms never do. Both are neutral. Does bonding depend on mass, nuclear charge, or electron arrangement? If arrangement, is hydrogen's one versus helium's two electrons about number or parity?

Second: carbon has four outer electrons. Diamond connects each carbon to four neighbors, graphite to three. Which traps electrons more tightly, which lets some roam, and which conducts?

Third: assemble ten quadrillion copper atoms and current flows. What moves, nuclei, copper ions, or electrons collectively owned by the whole metal? If the latter, why do they not all leak out?

\step{Hands-on Experiment}

\begin{experiment}[Does Pencil Lead Conduct?]
\textbf{Materials}: a pencil, preferably softer than 2B with more graphite; a 9~V battery or three series-connected 1.5~V cells; an inexpensive LED; salt; sugar; two cups of water; and wires, or twisted aluminum-foil strips if unavailable.

% BEGIN TRANSLATOR SAFETY NOTE
\textbf{Translator safety note:} Do not build the circuit as written: it omits a specified current-limiting component. Use an instructor-designed current-limited tester and insulated leads, never improvised foil connections or a direct battery short. Disconnect if anything warms. Batteries can overheat or leak when shorted; see \href{https://oem.energizer.com/english/dos-donts/}{Energizer battery-care guidance}.
% END TRANSLATOR SAFETY NOTE

\textbf{Step one}: sharpen both pencil ends to expose graphite. Connect battery, LED, and core in a loop, so current traverses core and LED before returning. LEDs are directional; reverse the legs if dark. It lights, showing graphite conduction. Replace the core with an eraser or plastic ruler and it goes out.

\textbf{Step two}: insert two separate wire ends into dry salt without touching each other, retaining battery and LED connections. No light: solid salt insulates.

% BEGIN TRANSLATOR SAFETY NOTE
\textbf{Translator safety note:} Do not power the saltwater arrangement below as a home experiment. Direct current can electrolyze brine, producing hazardous products including chlorine. Use a recorded demonstration or an instructor-approved conductivity tester; do not taste the solutions. See \href{https://edu.rsc.org/experiments/electrolysis-of-brine/735.article}{Royal Society of Chemistry brine-electrolysis safety guidance}.
% END TRANSLATOR SAFETY NOTE

\textbf{Step three}: dissolve salt in one cup and insert the separated wires. The LED lights. Replace with sugar water and it stays dark.

\textbf{Record and think}: graphite and salt are ordered crystals yet differ electrically; salt changes from insulating dry to conducting dissolved; equally clear saltwater and sugar water differ sharply. Does conduction belong to solid or liquid, atoms or something freed from them? Explain, including whether saltwater and graphite carry the same moving charges.
\end{experiment}

Our whole mystery hides here. Saltwater moves sodium and chloride ions, entire charged atoms; graphite moves electrons while atoms stay put. Identical conduction has different microscopic accounts; solid conduction is our main subject.

\step{The Old Model Runs into Trouble}

Attack with Coulomb's law and Chapter~52's atoms, and locate their failures.

First, electrostatics cannot choose partners. Neutral hydrogen atoms should cancel Coulomb forces, leaving only weak polarization attraction, too weak even to liquefy helium, certainly too weak for strong bonds. Worse, hydrogen binds hydrogen but neither hydrogen nor helium binds helium. Charge-only electrostatics sees identical neutral objects, with no selection mechanism. Classical physics cannot even justify molecules rather than perpetual atomic solitude.

Second, even granting attraction, why fixed numbers and directions? Hydrogen binds one hydrogen, never another; oxygen binds two hydrogens with a fixed 104.5-degree angle between three bonds. Why should attraction become saturated or directional? A directionless pull should welcome all comers equally. The old model has no clear reply.

Third, solid properties diverge inexplicably. Sticky hard balls predict similar crystals, not silver conducting a hundred quintillion times better than diamond, over twenty resistivity orders; graphite conducting hundreds of times better along layers than across; diamond's supreme hardness versus graphite's lubrication. Nor can they link hardness, transparency, and insulation in diamond. What conspiracy joins them?

All three expose one missing idea: electrons are not classical particles parked in atoms, but shaped, energetic quantum states as Chapter~52 explained. Approaching atoms reshuffle those states.

\step{Inventing New Concepts}

Invent four tools to fill the gap.

First, \textbf{molecular orbitals}. Overlapping clouds expose electrons to two nuclear potential valleys. Old atomic orbitals cannot simply coexist unchanged; they recombine into whole-molecule states. Two atomic orbitals always produce two new ones: a lower-energy \textbf{bonding orbital} and higher-energy \textbf{antibonding orbital}. Bonding density bridges nuclei and attracts both positive centers; antibonding density has a central gap and raises energy without binding. Like overlapping lamps whose waves can brighten or darken regions, probability waves reinforce and cancel. Unlike balls joined by a spring, an orbital is not a material tether but a molecular quantum distribution, not a route electrons run along.

Hydrogen and helium now differ precisely by the proposed parity. Hydrogen's two opposite-spin electrons share one discounted bonding room under Chapter~53's exclusion, lowering energy and stabilizing the molecule. Two helium atoms supply four electrons, forcing two into antibonding after bonding fills. Their added cost cancels, actually slightly exceeds, the first pair's savings, so helium molecules do not exist. Bonding means available cheaper quantum rooms, not sociability. Saturation follows: hydrogen's bonding orbital has two seats; a third atom would require antibonding occupancy and cost energy, so it is refused. Directionality, including water's angle, comes from orbital shapes; directional Chapter~52 p dumbbells point bonds toward partners.

Second, \textbf{ionic bonding}: complete electron relocation. Removing sodium's outer electron costs about \(5.1\)~eV; adding one to almost-full chlorine returns about \(3.6\)~eV affinity energy. Transfer alone costs \(1.5\)~eV, seemingly unfavorable. The payoff follows: opposite ions attract collectively. In salt each has six opposite nearest neighbors, so the whole crystal's electrostatic account matters. This \textbf{lattice energy}, about \(8\)~eV per pair, leaves over \(6\)~eV net gain. Not one sodium holding one chlorine, but a crystal-wide bulk purchase. Alternating checkerboard charges explain brittleness: sliding a layer half a square brings like charges face to face, reversing attraction into repulsion and cleaving salt. It also explains step three: fixed solid ions cannot carry current, but dissolved ions escape their cages and make the LED shine.

Third, \textbf{metallic bonding}, ultimate delocalization. Weakly bound outer orbitals of hundreds of millions of atoms recombine over the entire metal, not individual atoms or pairs. Electrons form a sea flowing through a positive-ion framework. Like a city's subway passengers, none belongs to one station. Unlike free gas dispersing randomly, collective positive attraction retains them. Leaving the surface costs work function energy, answering guess three. Nondirectional binding survives sliding atomic planes, permitting bending and wire drawing instead of salt-like fracture. Abundant nearly free carriers drift under fields and conduct electricity; mobile electrons transport heat too.

Fourth, \textbf{energy bands}: from two atoms to ten quadrillion. Two split one level into two, three into three, \(N\) into \(N\) closely spaced levels forming a continuous band. A cubic centimeter of copper contains roughly \(10^{23}\) atoms; levels crowd into a few eV with spacings around \(10^{-22}\)~eV, beyond instrumental resolution. Solid electrons occupy effectively continuous \textbf{bands}, separated possibly by forbidden \textbf{bandgaps} with no electron states. These unlock solid electrical properties; detailed accounting follows.

Before closing the toolbox, chemical bonds are not new forces. Electromagnetism is the only actor, performing a quantum script of indistinguishability, exclusion, delocalization, and state recombination. Bonds are collective electromagnetic performances under quantum rules. No new physics is required; everything is calculable in principle, though difficult enough to sustain an entire computational-chemistry profession.

\step{The Model in One Sentence}

\begin{oneline}
Atoms assemble by reshuffling electronic states into cheaper bonding and expensive antibonding rooms, filled lowest first under exclusion; favorable occupancy makes bonds. Ten quadrillion atoms crowd levels into bands: filled bands cannot conduct, half-filled ones permit flow, and gaps distinguish conductors, semiconductors, and insulators.
\end{oneline}

Like theater ticketing, bands are floors, equal-priced seats filled one person each, cheapest first. Partly filled floors allow movement; full ones require buying upstairs across a gap. Three limits: electrons are not choosing spectators, quantum rules specify occupancy; filled-band electrons are not stationary, their motions cancel, so nearby empty states determine conduction; and the seating chart belongs to the whole crystal, not separate independent atoms, the old model's mistake.

\step{The Mathematical Translator}

First graph whether bonding pays. Nuclear separation \(r\) sets total energy \(U(r)\). At large \(r\), independent atoms define \(U=0\). Decreasing \(r\) overlaps clouds, fills bonding orbitals, and lowers energy; as \(r\) shrinks further, nuclear repulsion and exclusion's expensive rooms dominate, raising it sharply. The minimum at \(r_0\) defines bond length, its depth below zero \(D\) bond energy.

\begin{figure}[htbp]
\centering
\begin{tikzpicture}[scale=1.0,>=Stealth]
  \draw[->] (0,0) -- (8.2,0) node[below left] {Separation \(r\)};
  \draw[->] (0,-3.2) -- (0,2.0) node[above left] {Energy \(U\)};
  \draw[thick,blue!60!black] (0.45,1.9)
    .. controls (0.75,-1.4) and (1.1,-2.55) .. (1.7,-2.62)
    .. controls (2.6,-2.7) and (3.6,-1.5) .. (5.2,-0.45)
    .. controls (6.4,-0.15) .. (8.0,-0.02);
  \draw[dashed] (1.7,-2.62) -- (1.7,0) node[above] {\(r_0\)};
  \draw[dashed] (0,-2.62) -- (1.7,-2.62);
  \draw[<->] (-0.12,-2.62) -- (-0.12,0) node[midway,left] {\(D\)};
  \node[blue!60!black] at (6.0,-1.85) {\shortstack{Bonded molecule's\\potential curve}};
  \draw[thick,red!60!black] (0.55,1.9)
    .. controls (1.2,0.6) and (2.2,0.25) .. (4.0,0.12)
    .. controls (6.0,0.04) .. (8.0,0.01);
  \node[red!60!black] at (4.6,1.1) {No bond, e.g. He--He};
\end{tikzpicture}
\end{figure}

Check intuitively. As \(r\to\infty\), \(U\to 0\), independent atoms; as \(r\to 0\), \(U\to+\infty\), enormous nuclear repulsion. At \(r=r_0\), zero slope means no force, equilibrium. Slightly smaller \(r\) than \(r_0\) gives negative slope and force toward increasing \(r\); slightly larger \(r\) than \(r_0\) gives positive slope and restoring attraction. The valley is stable. Helium's curve has no valley, only repulsion and a weak tail, no bond.

Real covalent bonds cost a few eV: H--H about \(4.5\)~eV, C--C \(3.6\)~eV; salt gains about \(6\)~eV per pair. At \(300\)~K, thermal energy per degree of freedom is roughly \(kT\approx0.026\)~eV. Bonds exceed room-temperature agitation by over a hundredfold, keeping tables, water, and bones stable rather than casually broken apart. Thousands of kelvin raise \(kT\) to tenths of an eV; energetic-tail molecules then break many bonds, underlying combustion and high-temperature metallurgy. The quiet eV-to-0.026eV ratio marks chemistry's stability boundary.

Second, count a salt grain. Take side \(0.5\)~mm and mass \(1\)~mg. At molar mass \(58.5\)~g/mol, it contains \(10^{-3}/58.5\approx1.7\times10^{-5}\)~mol, around \(10^{19}\) ion pairs. Neighbor separation \(0.28\)~nm and \(10^{19}\) pairs in a cube give roughly \((10^{19})^{1/3}\approx2\times10^{6}\) per edge, hence \(0.6\)~mm, matching the grain. Two million ions along each visible dimension connect macroscopic smallness to microscopic astronomical counts. Breaking crystals requires appreciable force despite eV bonds because ten quadrillion bonds break collectively.

Third, how continuous is a band? Copper's cubic centimeter has \(8.5\times10^{22}\) atoms, splitting an outer level into \(8.5\times10^{22}\) sublevels across a few eV. The \(10^{-22}\)~eV spacing is one-quintillionth of room-temperature \(kT\); even tiny thermal disturbances move electrons between neighbors. Band continuity is not a slogan but physically unresolvable spacing. Gaps are genuinely empty, without a single available state.

\begin{figure}[htbp]
\centering
\begin{tikzpicture}[scale=1.0,>=Stealth]
  % One atom: one level
  \draw[thick] (0,0) -- (1.2,0);
  \node[below] at (0.6,-0.1) {One atom};
  \node[right] at (1.3,0) {One level};
  % Two atoms: two levels
  \draw[thick] (3.2,0.25) -- (4.4,0.25);
  \draw[thick] (3.2,-0.25) -- (4.4,-0.25);
  \node[below] at (3.8,-0.35) {Two atoms};
  \node[right] at (4.5,0) {Two levels};
  % N atoms: a band
  \foreach \y in {-0.5,-0.4,-0.3,-0.2,-0.1,0,0.1,0.2,0.3,0.4,0.5}{
    \draw (6.6,\y) -- (7.8,\y);}
  \node[below] at (7.2,-0.6) {\(N\) atoms};
  \node[right] at (7.9,0) {A band};
  \draw[->,thick] (1.9,0) -- (2.9,0);
  \draw[->,thick] (5.0,0) -- (6.0,0);
\end{tikzpicture}
\end{figure}

Bands distinguish materials simply. Exclusion fills upward. A half-filled upper band, as in copper and silver, offers neighboring empty states; a small field shifts occupancy into directed current, a conductor. A full upper band separated from emptiness by a wide gap has no nearby room; opposing motions cancel, an insulator. Diamond's gap is about \(5.5\)~eV, beyond visible photons' roughly \(3\)~eV maximum, so light cannot promote electrons and passes through. Transparency and insulation are two sides of one account. Semiconductors are narrow-gap insulators: silicon's roughly \(1.1\)~eV gap admits a few thermally promoted electrons, intermediate conductivity extremely sensitive to heat and impurities. This sensitivity is a gift. Parts-per-million pentavalent phosphorus supplies spare electrons, n-type; trivalent boron leaves electron vacancies that neighbors fill, making oppositely moving positive holes, p-type. Humans now control conduction.

Join n and p regions in one silicon crystal for a diode. Forward voltage pushes carriers toward the interface, where electrons and holes recombine and replacements sustain current. Reverse voltage pulls carriers away, creating an empty depletion region and stopping current. No moving parts are needed for one-way passage. Check zero forward voltage: no source, no current. Huge reverse voltage eventually forces valence electrons across the gap, breakdown, showing every diode's limit. Add a control electrode to adjust a conducting channel's width by electric field and obtain a transistor: small control-voltage changes admit or stop much larger currents. Amplification and switching underpin the digital world; your phone industrially repeats this tens of billions of times.

\begin{mathbridge}
Bonds resemble springs with quantized energy. Approximate the potential valley as a parabola, the oscillation chapters' harmonic approximation:
\[
E_n=\Bigl(n+\tfrac12\Bigr)\hbar\omega,\qquad n=0,1,2,\dots
\]
Here \(\omega=\sqrt{k/\mu}\), with bond stiffness \(k\) and reduced atomic mass \(\mu\). At \(n=0\), \(E_0=\hbar\omega/2\neq0\), zero-point energy forbids complete vibrational rest even at absolute zero, since definite position and momentum violate uncertainty, independently of instruments. Large \(n\) makes spacing \(\hbar\omega\) small relative to \(E_n\), recovering a continuous classical spring. Typical \(\hbar\omega\) is tenths of an eV, infrared photon energy. Infrared promotes vibrational steps, enabling molecular identification and greenhouse absorption of terrestrial heat radiation.

Molecules rotate too. A diatomic molecule has \(I=\mu r_0^2\) about its center of mass and levels
\[
E_\ell=\frac{\hbar^2}{2I}\,\ell(\ell+1),\qquad \ell=0,1,2,\dots
\]
At \(\ell=0\), \(E_0=0\), true rotational rest is possible without fixing conjugate variables simultaneously, unlike vibration. Spacing grows linearly with \(\ell\); larger \(\ell\) gives taller steps, a subtle consequence of \(\ell(\ell+1)\) rather than \(\ell^2\). HCl rotational gaps are around \(10^{-3}\)~eV, microwave energies. Microwave ovens heat by water molecules absorbing rotational photons and redistributing energy through collisions. Infrared vibration and microwave rotation differ by two orders; scanning absorbed frequencies reveals bond strength, length, and mass, molecular spectroscopy's account.
\end{mathbridge}

\begin{challenge}
Why cannot full bands conduct? Beyond no empty neighbors, symmetry explains it. Periodic lattice potentials yield modulated plane-wave Bloch states labeled by crystal momentum \(\hbar k\). Full bands occupy all \(k\), with \(k\) and \(-k\) opposite velocities. Fields insufficient to cross the gap shift occupancy in \(k\) space but leave a fully symmetric filling, hence zero total current. Conduction means asymmetric \(k\)-space occupancy, not mere motion. Later, superconductivity pairs electrons into a different dissipation-free directed flow. Our lattice is a fixed background and independent electrons omit real-time correlations. Good for copper and silicon, this can fail badly for oxides including many high-temperature-superconductor parent compounds: Mott insulators, which simple bands predict metallic. Band theory is a model, not a law; its boundary opens correlated-electron physics.
\end{challenge}

\step{The Model's Limits}

Three bond types and bands need marked boundaries before you misapply them next chapter.

Molecular orbitals work especially well for delocalized hydrogen, oxygen, and benzene. Valence-bond theory emphasizes electrons shared between particular atoms. Both approximate the same quantum mechanics, agree where successful, and defer to fuller calculations when they differ. Molecular orbitals are mathematical models, not photographed clouds.

Real bonds rarely come pure. Water's covalent O--H bonds pull density toward oxygen, making polar ends and intermolecular hydrogen bonds. These tenths-of-an-eV bonds, roughly a tenth of covalent strength, make ice lighter: an open framework partially collapses on melting, densifying water. Hence floating ice, lakes freezing top-down, and winter survival below. Hydrogen bonds also support DNA's double-helix rungs. Still weaker van der Waals attraction acts between all neutral atoms and molecules, liquefying helium at extreme cold and adhering geckos to glass. Ionic/covalent/metallic categories map main roads, not every street.

Bands assume perfect periodic atomic arrays. Noncrystalline glass and plastic need qualified band language, though glass retains a photon-absorption threshold and transparency. Thermal vibrations disrupt periodicity and scatter electrons, raising metallic resistance. Some sufficiently cold metals pair electrons and lose resistance entirely, superconductivity, requiring added pairing physics discussed later. The challenge already warned of Mott insulators.

Our diode/transistor intuition omits details. Silicon forward current remains tiny below roughly \(0.7\)~V, then rises exponentially rather than switching abruptly; reverse leakage persists at nanoamp scales. Engineering is more detailed without reversing the qualitative picture. Transistors are not simply small-current controls for larger current: control-electrode fields alter channel conductivity, changing material state. Power comes entirely from the supply; amplification never violates energy accounting.

\step{Explaining the Opening Phenomenon}

Settle the three strange sibling pairs.

Diamond and graphite: four outer carbon electrons. Diamond's four orbitals form strong tetrahedral bonds with neighbors. They fill bonding states and lock a three-dimensional network in place, producing supreme hardness. Full bands and a \(5.5\)~eV gap prevent carriers; insufficient visible-photon energy gives transparency. Three properties, one story. Graphite bonds three neighbors in-plane; its fourth electron delocalizes across the sheet. Strong in-plane covalent networks, intrinsically stronger than diamond's, conduct; weak interlayer van der Waals forces permit sliding, writing, and lubrication. Poor interlayer electron transfer gives hundreds-fold conductivity anisotropy. Another intuitive trap: metals link electrical and thermal conduction, yet insulating diamond conducts heat among the best at room temperature, roughly five times copper. Rigid-lattice vibrations carry heat instead of electrons. Shared carriers tie metallic properties together, not a universal solid rule.

Salt and iron: collective ionic electrostatics turn repulsive when half-square sliding aligns like charges, causing brittle fracture. Nondirectional metallic electron seas refill gaps as planes slide, retaining tough bonds. Salt shatters; wire bends. Water dismantles salt lattices into conducting mobile ions; copper's electron sea conducts regardless of immersion.

Phone chips: silicon's narrow gap makes doping potent; impurities provide electrons or holes; n--p interfaces rectify; control electrodes create field-switched transistors. Billions of switch combinations build logic and everything behind this screen. The gray crystal's apparent decisions come from a roughly one-eV gap, not hidden machinery.

\step{Thought Experiments and Challenges}

\begin{thought}[At Extreme Pressure, Do Chemical Bonds Survive?]
Earthly bonds trade in a few eV, with separations set by electron-cloud equilibrium. Raise pressure: deep ocean \(10^{8}\)~Pa scarcely changes chemistry; Earth's center \(3\times10^{11}\)~Pa compresses atoms and metallizes some insulators, yet electrons retain shells. White-dwarf pressures exceed \(10^{22}\)~Pa, pushing nuclei much closer than ordinary atomic radii. Electrons no longer belong to individual nuclei; exclusion survives after atomic containers break, leaving a shared degenerate sea resisting gravity. This is Chapter~53's table-supporting wall holding a whole dead star. When do bonds, molecules, and solids cease to apply, gradually or at clear boundaries? Higher still, electrons enter protons to make neutron-star neutrons, dethroning even electromagnetism. Chemistry is a low-pressure cosmic privilege, readable from eV and pressure accounts.
\end{thought}

\begin{puzzles}
\item Hydrogen's bonding orbital holds two opposite-spin electrons. Can three atoms share a bonding-orbital chain to make stable neutral H\(_3\)? Reason from exclusion and bonding/antibonding counts. (Hint: three atomic orbitals make three molecular orbitals. Where does electron three sit, and do savings pay its rent? Real H\(_3^+\) exists but neutral H\(_3\) does not. Can your reasoning predict this?)
\item Why does a flowing negative electron sea not charge up a metal and tear it apart? (Hint: count positive framework charge and recall connections at wire ends. Current flows through rather than accumulating; circuits consume energy, not charge.)
\item Glass, diamond, and pure salt crystals insulate and transmit light. Give their common band explanation, then explain an opaque insulating ceramic bowl. (Hint: visible photons need not promote electrons, regardless of perfect atomic order. At nonuniform interfaces, what happens besides absorption? Scattering.)
\item Is pure silicon near absolute zero conducting or insulating? What after parts-per-million phosphorus doping? What does semiconductor's semi mean? (Hint: without thermal promotion full bands remain full and empty bands empty. Where does the extra donor electron reside? Semi does not mean half the conductivity, but how conductivity is controlled.)
\end{puzzles}

We have sewn atoms into molecules and solids with quantum rules and electromagnetism, no new forces. Next set this machine running: photons drive electrons between levels and bands, giving stimulated emission, lasers, LEDs, and atomic clocks. Knowing where the steps are, we will watch what dances on them.
